% TOP_PROBLEM: {"id": 30007407, "title": "The Voronoi pivotal-count amplitude under Cardy’s formula", "category_id": 19, "category": {"id": 19, "name": "probability", "display_name": "Probability", "description": "Problems involving probability theory, stochastic processes, and random structures.", "slug": "probability", "order_index": 19, "created_at": "2026-07-31T15:26:25.671Z"}, "status": "open"}
% FIRST_POSTED: 2026-10-06
% ATTEMPT_STATUS: solved
% OPENAI_CLAIM: {"schema_version":1,"source_repo":"https://github.com/openai/math","source_commit":"adc7f1241b42e322a6451854ab7e4b4c146bf78a","release_date":"2026-10-06","match":"full","resolution":"proved","independently_reviewed":false,"second_pass":"agreed","families":[{"family":"224","title":"Critical and quenched near-critical universality for Poisson–Voronoi percolation","subject":"Probability and statistical mechanics","manuscripts":[{"title":"A Pivotal Amplitude for Voronoi Percolation from Cardy's Formula","pdf_path":"preprints/A-Pivotal-Amplitude-for-Voronoi-Percolation-from-Cardys-Formula-October-5-2026/voronoi-pivotal-amplitude.pdf","readme_path":"preprints/A-Pivotal-Amplitude-for-Voronoi-Percolation-from-Cardys-Formula-October-5-2026/README.md","date":"2026-10-05","theorem_ref":"thm:main (sections/00-introduction.tex:35)","statement_sources":[{"path":"preprints/A-Pivotal-Amplitude-for-Voronoi-Percolation-from-Cardys-Formula-October-5-2026/build/sections/00-introduction.tex","line":35,"label":"thm:main"}]}],"lean":null,"reasoning_trace":null}]}
\documentclass[11pt]{article}
\usepackage[margin=25mm]{geometry}
\usepackage{fontspec}
\setmainfont{Latin Modern Roman}
\usepackage{amsmath,amssymb,mathtools}
\usepackage{unicode-math}
\setmathfont{Latin Modern Math}
\usepackage{enumitem}
\usepackage{xurl,hyperref,bookmark}
\hypersetup{colorlinks=true,urlcolor=blue}
\setcounter{secnumdepth}{0}
\setlength{\parindent}{0pt}
\setlength{\parskip}{5pt}
\setlength{\emergencystretch}{3em}
\begin{document}
\section*{The Voronoi pivotal-count amplitude under Cardy’s formula}

\subsection{OpenAI's claim}
OpenAI presents a proof in the selected manuscript. Assuming the Voronoi Cardy input without a convergence rate, the expected number of color pivotals is claimed to be asymptotic to $C\varepsilon^{-3/4}$ with $C>0$.

\subsection{Mathematical statement of the claimed result}
\paragraph{Result 1.} For full-plane Poisson--Voronoi cells at intensity $\varepsilon^{-2}$, independently colour each cell black with probability $1/2$. In a bounded Jordan domain $D$, let $a,b,c,d$ be distinct counterclockwise boundary marks; the crossing event requires a connected black subset of $\overline D$ meeting the closed arcs $ab$ and $cd$. Let $x\in(0,1)$ be the image of $b$ under the conformal map taking $(a,c,d)$ to $(0,1,\infty)$ in $\mathbb H$. Assume these scalar crossing probabilities have Cardy's limit for every bounded Jordan quadrilateral. Let $N_\varepsilon$ count nuclei whose colour flip changes the left-to-right crossing of $[0,1]^2$, including outside nuclei whose cells meet that square. There exists $c_V\in(0,\infty)$ with
\[
 \varepsilon^{3/4}\mathbb EN_\varepsilon\longrightarrow c_V,\qquad
 (\mathbb EN_\varepsilon)^{-1}\sim c_V^{-1}\varepsilon^{3/4}.
\]
Expectation averages geometry and colours. The hypothesis does not assume an interface limit or a convergence rate.

\subsection{Outline of the argument}
Four-interface transmission through annuli is decomposed spectrally to separate the conformal information at each end. Uniform signed-state transfer fixes the amplitude, and an intensity derivative connects the transmission calculation to the expected pivotal count.

\subsection{Scope relative to this problem}
Assuming the Voronoi Cardy input without a convergence rate, the expected number of color pivotals is claimed to be asymptotic to $C\varepsilon^{-3/4}$ with $C>0$. The question retains these same parameters and conventions. Within these assumptions, this is a claimed affirmative answer to the stated question.

\subsection{Formal verification}
Family 224 has no Lean scope page in this release. Its selected result is not presented as formalized. This site has not executed the Lean code. The scope descriptions do not certify the remaining manuscript arguments.

\subsection{Status caveat}
OpenAI attributes these manuscripts to work by an internal model. Its release cautions that arguments without formalization can contain mistakes. This catalogue has not independently checked the proofs; the labels report the mathematical claims attributed to OpenAI.

\subsection{Sources}
\begin{itemize}
\item OpenAI, \emph{A Pivotal Amplitude for Voronoi Percolation from Cardy's Formula}, 2026-10-05; citation key \texttt{OAI:A-Pivotal-Amplitude-for-Voronoi-Percolation-from-Cardys-Formula-October-5-2026}. \href{https://github.com/openai/math/blob/adc7f1241b42e322a6451854ab7e4b4c146bf78a/preprints/A-Pivotal-Amplitude-for-Voronoi-Percolation-from-Cardys-Formula-October-5-2026/voronoi-pivotal-amplitude.pdf}{PDF (GitHub)}; \href{https://github.com/openai/math/raw/adc7f1241b42e322a6451854ab7e4b4c146bf78a/preprints/A-Pivotal-Amplitude-for-Voronoi-Percolation-from-Cardys-Formula-October-5-2026/voronoi-pivotal-amplitude.pdf}{PDF (direct)}; \href{https://github.com/openai/math/blob/adc7f1241b42e322a6451854ab7e4b4c146bf78a/preprints/A-Pivotal-Amplitude-for-Voronoi-Percolation-from-Cardys-Formula-October-5-2026/README.md}{Manuscript page and citation}; \href{https://github.com/openai/math/blob/main/preprints/A-Pivotal-Amplitude-for-Voronoi-Percolation-from-Cardys-Formula-October-5-2026/voronoi-pivotal-amplitude.pdf}{Latest version}; Statement source: \href{https://github.com/openai/math/blob/adc7f1241b42e322a6451854ab7e4b4c146bf78a/preprints/A-Pivotal-Amplitude-for-Voronoi-Percolation-from-Cardys-Formula-October-5-2026/build/sections/00-introduction.tex\#L35}{thm:main}.
\item \href{https://github.com/openai/math/blob/adc7f1241b42e322a6451854ab7e4b4c146bf78a/CONTENTS.md}{OpenAI catalogue entry}.
\end{itemize}
\end{document}
